\section{Reproducibility Package}

Reproducibility is treated as part of the contribution rather than as an afterthought. The submission package is organized so that a reviewer or future researcher can reconstruct the manuscript, the deployment logic, and the validation workflow with minimal hidden assumptions.

The package includes:
\begin{itemize}
    \item source code for the platform services, controllers, and supporting components;
    \item deployment manifests and image build logic for the K3s-based setup;
    \item the complete LaTeX manuscript, bibliography database, and submission support material;
    \item protocol-level scripts and operational commands used to exercise enrollment, heartbeat supervision, and deployment validation;
    \item the measurement harness output \texttt{experiments/scalability\_metrics\_1778929086.json}, which stores all raw trial records and aggregate statistics used in the manuscript;
    \item the figure-generation script \texttt{paper\_ieee\_inginf05/generate\_figures.py}, which regenerates the plots embedded in the paper from the stored JSON artifact;
    \item documentation describing the expected environment and the sequence required to reproduce the experiments.
\end{itemize}

The artifact can be executed on a commodity Linux host with Docker and K3s, which keeps the barrier to reproduction relatively low while still preserving realistic cloud-native execution semantics. For a final journal submission, the package should be strengthened further by pinning a stable repository tag, recording the exact commit hash, and exporting immutable artifact identifiers such as release archives or checksums. These additions are straightforward but important because a paper's long-term reproducibility depends on freezing the evaluated implementation state.

A second reproducibility consideration concerns observational evidence. The present study relies on gateway logs, API-server statistics, Kubernetes resource snapshots, raw experiment JSON, and generated manuscript assets. This diversity of evidence is beneficial because it allows independent confirmation of the same event sequence through multiple viewpoints. For example, a successful deployment is visible both as a device-side runtime acknowledgment exported through the API server and as a northbound application record. Such redundancy is especially useful during peer review, where reviewers often need confidence that reported results are not artifacts of a single reporting channel.

To regenerate the figures included in this manuscript, a reviewer can execute:
\begin{verbatim}
cd paper_ieee_inginf05
../.venv-paper/bin/python generate_figures.py
latexmk -pdf main.tex
\end{verbatim}
The resulting images are written under \texttt{paper\_ieee\_inginf05/figures/} and are directly consumed by the LaTeX source.
